\documentclass[12pt,a4paper]{ctexart}

% ==== 字体与编码 ====
\usepackage{fontspec}
\setmainfont{Times New Roman}
\setsansfont{Arial}
\setmonofont{Consolas}
\setCJKmainfont{SimSun}[AutoFakeBold=2.5]
\setCJKsansfont{SimHei}
\setCJKmonofont{KaiTi}

% ==== 数学 ====
\usepackage{amsmath,amssymb,amsthm}
\usepackage{mathtools}

% ==== 页面布局 ====
\usepackage[top=2.5cm,bottom=2.5cm,left=2.8cm,right=2.8cm]{geometry}

% ==== 超链接 ====
\usepackage[colorlinks=true,linkcolor=blue,citecolor=blue,urlcolor=blue]{hyperref}

% ==== 交换图 ====
\usepackage{tikz-cd}

% ==== 列表 ====
\usepackage{enumitem}

% ==== 表格线 ====
\usepackage{booktabs}

% ==== 定理环境 ====
\theoremstyle{definition}
\newtheorem{definition}{定义}[section]
\newtheorem{example}[definition]{例}
\newtheorem{remark}[definition]{注记}
\newtheorem{proposition}[definition]{命题}
\newtheorem{theorem}[definition]{定理}
\newtheorem{corollary}[definition]{推论}
\newtheorem{lemma}[definition]{引理}
\newtheorem{construction}[definition]{构造}

% ==== 自定义命令 ====
\newcommand{\cat}[1]{\mathbf{#1}}
\newcommand{\Set}{\cat{Set}}
\newcommand{\Vect}{\cat{Vect}}
\newcommand{\Cat}{\cat{Cat}}
\newcommand{\Para}{\cat{Para}}
\newcommand{\id}{\mathrm{id}}
\newcommand{\Monad}{\mathrm{Mnd}}
\newcommand{\Alg}{\mathrm{Alg}}
\newcommand{\Coalg}{\mathrm{Coalg}}
\newcommand{\LaxAlg}{\mathrm{Lax}\text{-}\Alg}
\newcommand{\FreeMnd}{\mathrm{FreeMnd}}
\newcommand{\Stream}{\mathrm{Stream}}
\newcommand{\Moore}{\mathrm{Moore}}
\newcommand{\Mealy}{\mathrm{Mealy}}
\newcommand{\List}{\mathrm{List}}
\newcommand{\Tree}{\mathrm{Tree}}
\newcommand{\ob}[1]{\mathrm{Ob}(#1)}
\newcommand{\mor}[1]{\mathrm{Mor}(#1)}
\newcommand{\colim}{\mathrm{colim}}
\newcommand{\limite}{\mathrm{lim}}
\newcommand{\fold}{\mathrm{fold}}
\newcommand{\cell}{\mathrm{cell}}
\newcommand{\unfold}{\mathrm{unfold}}
\newcommand{\Endo}{\mathrm{Endo}}

\title{\textbf{范畴深度学习：所有架构的代数理论}\\[0.3em]
\large ——完整论文教程，含严谨数学补充}
\author{基于 Bruno Gavranović, Paul Lessard, Andrew Dudzik 等 (ICML 2024)\\ 整理与讲解}
\date{}

\begin{document}
\maketitle

\tableofcontents
\newpage

% ==================================================================
\section{引言：范畴论与深度学习的交汇}
% ==================================================================

\subsection{论文原旨}

\begin{quote}
\itshape
It is our position that constructing a guiding framework for all of deep learning requires robustly bridging the top-down and bottom-up approaches to neural network specification with a unifying mathematical theory, and that the concepts for this bridging should be coming from computer science. Moreover, such a framework must generalise both group theory and functional programming---and a natural candidate for achieving this is category theory.
\end{quote}

\textbf{翻译：}

我们的立场是：为所有深度学习构建一个指导性框架，需要将神经网络规范的自顶向下和自底向上两种方法用统一的数学理论稳健地桥接起来，而这一桥接概念应该来自计算机科学。此外，这一框架必须同时推广群论和函数式编程——实现这一目标的自然候选者是范畴论。

\textbf{讲解：}

这篇论文的核心贡献在于提出：范畴论中的\textbf{单子（monad）}及其代数，取值于参数化映射的\textbf{2-范畴 Para}，能够统一地刻画：
\begin{itemize}[nosep]
    \item 几何深度学习中的\textbf{等变性}约束（群作用单子的代数同态）；
    \item 函数式编程中的\textbf{归纳/余归纳数据结构}（列表、树的自函子代数）；
    \item 自动机理论中的\textbf{状态机}（Mealy、Moore 机器的自函子余代数）；
    \item 神经网络的\textbf{权重共享}（通过 Para 2-范畴中的余交换余幺半群结构）。
\end{itemize}

\subsection{阅读指南}

本教程按论文的自然顺序组织。每一节先给出论文核心原文的翻译，然后用更严谨的数学语言补充定义、定理和交换图。特别在作者写简略的地方——如列表作为代数范畴始对象的构造、通过自函子迭代取余极限得到自由单子的超穷构造——我们将给出完整推导。

\newpage

% ==================================================================
\section{范畴论基础}
% ==================================================================

论文第 1.3.1 节回顾了范畴论的三个核心概念：范畴、函子、自然变换。下面给出严谨定义。

\begin{definition}[范畴]\label{def:category}
一个\textbf{范畴} $\cat C$ 由以下数据组成：
\begin{enumerate}[label=(\arabic*),nosep]
    \item 一个\textbf{对象}的类 $\ob{\cat C}$；
    \item 对任意 $A,B\in\ob{\cat C}$，一个\textbf{态射}的类 $\cat C(A,B)$；
    \item 对任意 $A\in\ob{\cat C}$，一个\textbf{恒等态射} $\id_A\in\cat C(A,A)$；
    \item 对任意 $A,B,C\in\ob{\cat C}$，一个\textbf{复合}映射
    \[\circ:\cat C(B,C)\times\cat C(A,B)\longrightarrow\cat C(A,C),\quad (g,f)\mapsto g\circ f.\]
\end{enumerate}
满足：
\begin{enumerate}[label=(\roman*),nosep]
    \item \textbf{单位律}：$\id_B\circ f=f=f\circ\id_A$，$\forall f\in\cat C(A,B)$；
    \item \textbf{结合律}：$h\circ(g\circ f)=(h\circ g)\circ f$，$\forall f\in\cat C(A,B),\,g\in\cat C(B,C),\,h\in\cat C(C,D)$。
\end{enumerate}
\end{definition}

\begin{example}[集合范畴 $\Set$]\label{ex:set}
$\Set$ 的对象是所有集合，态射是集合之间的映射，复合是函数的复合，恒等态射是恒等函数。
\end{example}

\begin{example}[群和幺半群作为范畴]\label{ex:group}
设 $G$ 是一个群。定义范畴 $\mathbf{B}G$：
\begin{itemize}[nosep]
    \item 只有一个对象 $*$；
    \item 态射 $\mathbf{B}G(*,*)$ 的元素对应于 $G$ 的元素；
    \item 复合对应于 $G$ 的乘法：$g\circ h:=gh$；
    \item 恒等态射 $\id_*$ 对应于 $G$ 的单位元 $e$。
\end{itemize}
$G$ 是群当且仅当所有态射都是同构（即可逆）。更一般地，只有一个对象的范畴（态射不必可逆）精确对应于\textbf{幺半群}（monoid）。
\end{example}

\begin{definition}[函子]\label{def:functor}
设 $\cat C,\cat D$ 为范畴。一个\textbf{函子} $F:\cat C\to\cat D$ 包括：
\begin{enumerate}[label=(\arabic*),nosep]
    \item 对象映射：$A\mapsto F(A)$；
    \item 态射映射：$f\mapsto F(f)$，且 $F(f)\in\cat D(F(A),F(B))$ 当 $f\in\cat C(A,B)$；
\end{enumerate}
满足 $F(\id_A)=\id_{F(A)}$ 和 $F(g\circ f)=F(g)\circ F(f)$。
\end{definition}

若 $\cat C=\cat D$，则称 $F$ 为\textbf{自函子}（endofunctor）。

\begin{definition}[自然变换]\label{def:natural}
设 $F,G:\cat C\to\cat D$ 为两个函子。一个\textbf{自然变换} $\alpha:F\Rightarrow G$ 由一族 $\cat D$ 中的态射
\[\{\alpha_X:F(X)\to G(X)\}_{X\in\ob{\cat C}}\]
（称为 $\alpha$ 在 $X$ 处的\textbf{分量}）组成，使得对任意 $\cat C$ 中的态射 $f:X\to Y$，下图交换：
\begin{equation}\label{eq:nat-square}
\begin{tikzcd}
F(X) \arrow[r,"F(f)"] \arrow[d,"\alpha_X"'] & F(Y) \arrow[d,"\alpha_Y"] \\
G(X) \arrow[r,"G(f)"'] & G(Y)
\end{tikzcd}
\end{equation}
即 $\alpha_Y\circ F(f)=G(f)\circ\alpha_X$。若每个分量 $\alpha_X$ 都是同构，则称 $\alpha$ 为\textbf{自然同构}。
\end{definition}

这三个概念——范畴、函子、自然变换——构成了范畴论的"三位一体"，也是论文后续一切构造的基础。

\newpage

% ==================================================================
\section{单子与等变性：几何深度学习的范畴推广}
% ==================================================================

\subsection{原文翻译}

\begin{quote}
\itshape
We will define a powerful notion (monad algebra homomorphism) and demonstrate that the special case of monads induced by group actions is sufficient to describe geometric deep learning. Generalising from monads and their algebras to arbitrary endofunctors and their algebras, we will find that our theory can express functions that process structured data from computer science (e.g.\ lists and trees) and behave in stateful ways like automata.
\end{quote}

\textbf{翻译：}

我们将定义一个强有力的概念——单子代数同态，并证明群作用诱导的单子这一特例足以描述几何深度学习。从单子及其代数推广到任意自函子及其代数后，我们的理论可以表达处理计算机科学中结构化数据（如列表和树）的函数，以及像自动机那样具有状态行为的函数。

\subsection{单子及其代数}

\begin{definition}[单子]\label{def:monad}
设 $\cat C$ 是一个范畴。$\cat C$ 上的一个\textbf{单子}（monad）是一个三元组 $(M,\eta,\mu)$，其中
\begin{itemize}[nosep]
    \item $M:\cat C\to\cat C$ 是一个自函子；
    \item $\eta:\id_{\cat C}\Rightarrow M$ 是自然变换，称为\textbf{单位}（unit）；
    \item $\mu:M\circ M\Rightarrow M$ 是自然变换，称为\textbf{乘法}（multiplication）；
\end{itemize}
使得下图交换：
\[
\begin{tikzcd}[column sep=large, row sep=large]
M^3(X) \arrow[r,"M\mu_X"] \arrow[d,"\mu_{MX}"'] & M^2(X) \arrow[d,"\mu_X"] \\
M^2(X) \arrow[r,"\mu_X"'] & M(X)
\end{tikzcd}
\qquad
\begin{tikzcd}[column sep=large, row sep=large]
M(X) \arrow[r,"M\eta_X"] \arrow[dr,equal] & M^2(X) \arrow[d,"\mu_X"] & M(X) \arrow[l,"\eta_{MX}"'] \arrow[dl,equal] \\
& M(X) &
\end{tikzcd}
\]
即 $\mu_X\circ M\mu_X=\mu_X\circ\mu_{MX}$（结合律），以及 $\mu_X\circ M\eta_X=\id_{MX}=\mu_X\circ\eta_{MX}$（单位律）。
\end{definition}

\begin{example}[群作用单子]\label{ex:group-monad}
设 $G$ 是一个群。定义 $\Set$ 上的单子 $(G\times-,\eta,\mu)$：
\begin{itemize}[nosep]
    \item 自函子 $G\times-:\Set\to\Set$，将集合 $X$ 映为 $G\times X$，将映射 $f:X\to Y$ 映为 $\id_G\times f:G\times X\to G\times Y$；
    \item $\eta_X:X\to G\times X$，$x\mapsto(e,x)$，其中 $e$ 是群单位元；
    \item $\mu_X:G\times G\times X\to G\times X$，$(g,h,x)\mapsto(gh,x)$。
\end{itemize}
\textbf{验证单子公理}：
\begin{itemize}
    \item 结合律：$\mu_X(\id_G\times\mu_X)(g,(h,(k,x)))=\mu_X(g,hk,x)=(g(hk),x)=\mu_X(gh,k,x)=\mu_X(\mu_G\times\id_X)(g,h,k,x)$，成立因为群的结合律。
    \item 单位律：$\mu_X(\eta_{G\times X}(g,x))=\mu_X(e,g,x)=(eg,x)=(g,x)$，等等。
\end{itemize}
\end{example}

\begin{definition}[单子的代数]\label{def:monad-algebra}
设 $(M,\eta,\mu)$ 为 $\cat C$ 上的单子。一个\textbf{$M$-代数}是一个对 $(A,a)$，其中 $A\in\ob{\cat C}$ 是\textbf{承载对象}，$a:M(A)\to A$ 是 $\cat C$ 中的态射（称为\textbf{结构映射}），使得下图交换：
\[
\begin{tikzcd}[column sep=large, row sep=large]
A \arrow[r,"\eta_A"] \arrow[dr,"\id_A"'] & M(A) \arrow[d,"a"] & M^2(A) \arrow[l,"\mu_A"'] \arrow[d,"M(a)"] \arrow[r,"M(a)"] & M(A) \arrow[d,"a"] \\
& A & M(A) \arrow[r,"a"'] & A
\end{tikzcd}
\]
即 $a\circ\eta_A=\id_A$ 且 $a\circ\mu_A=a\circ M(a)$。
\end{definition}

\begin{definition}[$M$-代数同态]\label{def:monad-alg-hom}
设 $(M,\mu,\eta)$ 为单子，$(A,a)$ 和 $(B,b)$ 为 $M$-代数。一个\textbf{$M$-代数同态} $f:(A,a)\to(B,b)$ 是 $\cat C$ 中的态射 $f:A\to B$，使得下图交换：
\begin{equation}\label{eq:alg-hom}
\begin{tikzcd}
M(A) \arrow[r,"M(f)"] \arrow[d,"a"'] & M(B) \arrow[d,"b"] \\
A \arrow[r,"f"'] & B
\end{tikzcd}
\end{equation}
\end{definition}

\subsection{等变映射作为单子代数同态}

\begin{example}[等变映射]\label{ex:equivariant}
取 $G$ 为一个群，$(G\times-,\eta,\mu)$ 为群作用单子。设 $X,Y$ 是两个 $G$-集（即带有左 $G$-作用 $\alpha:G\times X\to X$，$\beta:G\times Y\to Y$），则它们都是 $G\times-$ 的代数。此时代数同态 $f:X\to Y$ 需要满足：
\[
\begin{tikzcd}
G\times X \arrow[r,"\id_G\times f"] \arrow[d,"\alpha"'] & G\times Y \arrow[d,"\beta"] \\
X \arrow[r,"f"'] & Y
\end{tikzcd}
\]
逐元素写开即为：$f(g\cdot x)=g\cdot f(x)$。这正是\textbf{$G$-等变映射}的定义。

特别地，取 $G=\mathbb Z_w\times\mathbb Z_h$（带加法的平移群），作用在图像数据 $\mathbb R^{\mathbb Z_w\times\mathbb Z_h}$ 上：
\[((i',j')\rhd x)(i,j)=x(i-i',j-j'),\]
则等变条件 $f(g\rhd x)=g\rhd f(x)$ 导出平移等变的线性层——即卷积层。
\end{example}

\textbf{补充讲解：} 论文此处的核心洞察是：几何深度学习（GDL）中所有的等变约束都可以统一地表述为\textbf{群作用单子的代数同态}。这比 GDL 的原始框架更为一般，因为单子不限于群——任何幺半群的作用都可以用单子 $M\times-$ 来编码，甚至更一般的代数理论（如 Lawvere 理论）也可以通过单子来刻画。

\begin{remark}[单子代数与函子范畴的等价]\label{rem:monad-functor}
当单子为 $M\times-$（$M$ 为幺半群）时，$M$-代数范畴等价于函子范畴 $[\mathbf{B}M,\Set]$，即从独对象范畴到 $\Set$ 的函子范畴。这是因为 $M$-作用等同于函子 $\mathbf{B}M\to\Set$，而代数同态等同于自然变换。这一观察将我们的定义与 de Haan et al. (2020) 的函子视角统一起来。
\end{remark}

\newpage

% ==================================================================
\section{自函子的（余）代数：结构化数据与自动机}
% ==================================================================

\subsection{原文翻译}

\begin{quote}
\itshape
Geometric deep learning, while elegant, is fundamentally constrained by the axioms of group theory. Monads and their algebras, however, are naturally generalised beyond group actions. Here we show how, by studying (co)algebras of arbitrary endofunctors, we can rediscover standard computer science constructs like lists, trees and automata.
\end{quote}

\textbf{翻译：}

几何深度学习虽然优雅，但根本上受限于群论的公理。然而单子及其代数自然地推广到群作用之外。我们在此展示，通过研究任意自函子的（余）代数，我们能够重新发现计算机科学中的标准构造，如列表、树和自动机。

\subsection{自函子的代数}

\begin{definition}[自函子的代数]\label{def:endo-alg}
设 $\cat C$ 为范畴，$F:\cat C\to\cat C$ 为一个自函子。$F$ 的一个\textbf{代数}是一个对 $(A,a)$，其中 $A\in\ob{\cat C}$，$a:F(A)\to A$ 是 $\cat C$ 中的态射。
\end{definition}

注意：与单子代数（定义 \ref{def:monad-algebra}）不同，这里没有额外的等式约束，因为 $F$ 本身不带有 $\eta$ 和 $\mu$ 这样的结构。

\begin{definition}[自函子代数的同态]\label{def:endo-alg-hom}
设 $(A,a)$ 和 $(B,b)$ 为 $F$-代数。一个\textbf{$F$-代数同态} $f:(A,a)\to(B,b)$ 是态射 $f:A\to B$，使得下图交换：
\[
\begin{tikzcd}
F(A) \arrow[r,"F(f)"] \arrow[d,"a"'] & F(B) \arrow[d,"b"] \\
A \arrow[r,"f"'] & B
\end{tikzcd}
\]
所有 $F$-代数及它们之间的同态构成一个范畴，记作 $\Alg(F)$。
\end{definition}

\subsection{列表作为始代数}

论文中的例 2.9 和注 2.13 讨论了列表作为自函子代数的始对象。这里作者写得比较简略，下面我们给出完整的数学构造。

\begin{example}[列表]\label{ex:list}
固定集合 $A$。考虑自函子
\[F_A:\Set\longrightarrow\Set,\quad F_A(X)=1+A\times X,\]
其中 $1=\{\bullet\}$ 是单点集（终端对象），$+$ 是余积（不相交并），$\times$ 是积。

\textbf{列表集合} $\List(A)$ 由所有有限序列 $[a_1,a_2,\dots,a_n]$ 组成（$n\ge 0$，$a_i\in A$）。定义结构映射
\[[\mathrm{Nil},\mathrm{Cons}]:1+A\times\List(A)\longrightarrow\List(A)\]
为：
\[\mathrm{Nil}(\bullet)=[\,],\qquad \mathrm{Cons}(a,\ell)=a::\ell\;\;(\text{将 }a\text{ 加到列表 }\ell\text{ 的前面}).\]
则 $(\List(A),[\mathrm{Nil},\mathrm{Cons}])$ 是 $F_A$ 的一个代数。
\end{example}

\begin{theorem}[列表是 $F_A$-代数的始对象]\label{thm:list-initial}
对任意 $F_A$-代数 $(X,[r_0,r_1])$（其中 $r_0:1\to X$，$r_1:A\times X\to X$），存在\textbf{唯一}的代数同态 $\fold:\List(A)\to X$，使得下图交换：
\[
\begin{tikzcd}
1+A\times\List(A) \arrow[r,"\id_1+\id_A\times\fold"] \arrow[d,"{[\mathrm{Nil},\mathrm{Cons}]}"'] & 1+A\times X \arrow[d,"{[r_0,r_1]}"] \\
\List(A) \arrow[r,"\fold"'] & X
\end{tikzcd}
\]
\end{theorem}

\begin{proof}[构造性证明（结构递归）]
定义 $\fold$ 通过列表上的递归：
\begin{align*}
\fold([\,]) &= r_0(\bullet),\\
\fold(a::\ell) &= r_1(a,\fold(\ell)).
\end{align*}
\textbf{唯一性：} 若 $g:\List(A)\to X$ 也是代数同态，则交换条件强制
\[g([\,])=r_0(\bullet),\quad g(a::\ell)=r_1(a,g(\ell)),\]
所以由列表归纳，$g=\fold$。
\end{proof}

\textbf{注：} 始代数性质的要点是：\textbf{任何从列表到某个类型的函数，如果结构递归地定义，则存在且唯一}。这在函数式编程中称为 \texttt{fold}（折叠），是现代 RNN 展开的范畴论基础。

\subsection{列表作为自由单子的构造（超穷迭代）}

论文在注 2.13 和附录 B.2 中提到"代数自由单子"（algebraically free monad）是通过不断迭代自函子直到稳定来构造的。这里我们补充完整的数学细节。

\begin{construction}[从自函子到自由单子的超穷构造]\label{cons:free-monad}
设 $\cat C$ 是有余极限的范畴，$F:\cat C\to\cat C$ 为自函子。定义一列对象 $\{F^\kappa(\emptyset)\}_{\kappa}$ 如下：
\begin{itemize}[nosep]
    \item $F^0(\emptyset)=\emptyset$（始对象）；
    \item $F^{\kappa+1}(\emptyset)=F(F^\kappa(\emptyset))$；
    \item 若 $\lambda$ 是极限序数，则 $F^\lambda(\emptyset)=\displaystyle\colim_{\kappa<\lambda}F^\kappa(\emptyset)$。
\end{itemize}
则对充分大的正则基数 $\kappa$（使得 $F$ 保持 $\kappa$-滤过余极限时），$F^\kappa(\emptyset)$ 稳定，即
\[\FreeMnd(F):=F^\kappa(\emptyset)\;\text{ 满足 }\;F(\FreeMnd(F))\cong\FreeMnd(F).\]
这是 $F$ 生成的\textbf{自由单子}的承载对象，且 $(\FreeMnd(F),F^\kappa)$ 范畴等价于 $F$-代数的范畴 $\Alg(F)$。
\end{construction}

\begin{example}[列表作为自由单子的特例]
对于 $F_A(X)=1+A\times X$，从空集 $\emptyset$ 开始迭代：
\begin{align*}
F_A^0(\emptyset) &= \emptyset,\\
F_A^1(\emptyset) &= 1+A\times\emptyset \cong 1,\\
F_A^2(\emptyset) &= 1+A\times 1 \cong 1+A,\\
F_A^3(\emptyset) &= 1+A\times(1+A) \cong 1+A+A^2,\\
&\;\;\vdots\\
F_A^\omega(\emptyset) &= \colim_{n<\omega} F_A^n(\emptyset) \cong \coprod_{n=0}^\infty A^n = \List(A).
\end{align*}
此时 $F_A(\List(A))\cong\List(A)$（在 $\omega$ 处稳定），且 $(\List(A),[\mathrm{Nil},\mathrm{Cons}])$ 上的单子结构为：
\begin{itemize}[nosep]
    \item $\eta_X:X\to\List(A)\times X$（通过空列表嵌入），
    \item $\mu_X:\List(A)\times\List(A)\times X\to\List(A)\times X$（列表拼接）。
\end{itemize}
\end{example}

这解释了为什么列表不仅是 $F_A$ 的始代数，也是 $F_A$ 生成的\textbf{自由单子}的承载集。在后续的 Part 3 中，这个构造将被提升到 2-范畴 Para 中，用以自动生成循环和递归神经网络的展开。

\subsection{二叉树的代数}

\begin{example}[二叉树]\label{ex:tree}
固定集合 $A$。考虑自函子
\[G_A(X)=A+X^2.\]
$G_A$ 的始代数是 $(\Tree(A),[\mathrm{Leaf},\mathrm{Node}])$，其中
\begin{itemize}[nosep]
    \item $\Tree(A)$ 是所有叶节点标记为 $A$ 的有限二叉树的集合；
    \item $\mathrm{Leaf}:A\to\Tree(A)$ 构造叶节点；
    \item $\mathrm{Node}:\Tree(A)^2\to\Tree(A)$ 构造内部节点。
\end{itemize}
对应的折叠（fold）函数递归定义为：
\begin{align*}
\fold(\mathrm{Leaf}(a)) &= r_0(a),\\
\fold(\mathrm{Node}(l,r)) &= r_1(\fold(l),\fold(r)).
\end{align*}
\end{example}

\subsection{自函子的余代数与自动机}

\begin{definition}[自函子的余代数]\label{def:endo-coalg}
设 $F:\cat C\to\cat C$ 为自函子。$F$ 的一个\textbf{余代数}是一个对 $(C,c)$，其中 $C\in\ob{\cat C}$，$c:C\to F(C)$ 是态射（\textbf{注意方向：与代数相反}）。
\end{definition}

余代数的同态定义对称地将箭头反向：
\[
\begin{tikzcd}
C \arrow[r,"c"] \arrow[d,"f"'] & F(C) \arrow[d,"F(f)"] \\
D \arrow[r,"d"'] & F(D)
\end{tikzcd}
\]

\textbf{直观理解：} 代数 $a:F(A)\to A$ 是一种"归纳构造"——从 $F$-结构\textbf{折叠}出一个结果（计算必定终止）。余代数 $c:C\to F(C)$ 是一种"余归纳观察"——从一个状态\textbf{展开}出下一步的行为和新的状态（计算可以无限进行）。

\begin{example}[Mealy 机器]\label{ex:mealy}
设 $O$（输出集）和 $I$（输入集）为集合。考虑自函子
\[H_{O,I}(X)=I\to O\times X.\]
$H_{O,I}$ 的终余代数是 $\Mealy_{O,I}$，即所有 Mealy 机器的集合。Mealy 机器的行为是：给定输入 $i\in I$，产生输出 $o\in O$ 并转移到新的 Mealy 机器状态。

用 Haskell 风格的数据类型表示：
\begin{verbatim}
data Mealy o i = MkMealy { next :: i -> (o, Mealy o i) }
\end{verbatim}
其结构映射 $\mathrm{next}:\Mealy_{O,I}\to(I\to O\times\Mealy_{O,I})$ 正是余代数结构。

从任意 $H_{O,I}$-余代数 $(X,n)$ 到 $(\Mealy_{O,I},\mathrm{next})$ 的终余代数同态（即\textbf{展开}，unfold）为：
\begin{verbatim}
unfold :: (x -> i -> (o, x)) -> x -> Mealy o i
unfold n x = MkMealy (\i -> let (o', x') = n x i
                             in (o', unfold n x'))
\end{verbatim}
\end{example}

\begin{example}[流（Stream）]\label{ex:stream}
取 $F_O(X)=O\times X$。$F_O$ 的终余代数是 $\Stream(O)$，即所有输出为 $O$ 的无限序列（同构于 $\mathbb N\to O$）。结构映射为
\[\langle\mathrm{head},\mathrm{tail}\rangle:\Stream(O)\longrightarrow O\times\Stream(O),\]
$\mathrm{head}$ 取第一个元素，$\mathrm{tail}$ 取去掉第一个元素后的剩余流。
\end{example}

\begin{example}[Moore 机器]\label{ex:moore}
取 $F_{O,I}(X)=O\times(I\to X)$。终余代数是 $\Moore_{O,I}$，其行为是：先输出，再接收输入，然后转移到下一个 Moore 机器状态。与 Mealy 机器的区别在于：Moore 机器的输出\textbf{不依赖于当前输入}。
\end{example}

\subsection{代数同态与结构递归}

论文例 2.12 讨论了列表上的 fold 作为代数同态。这个对应是整篇论文最重要的技术枢纽，下面给出严谨的展开。

\begin{proposition}[fold 作为始代数同态]\label{prop:fold}
设 $F_A(X)=1+A\times X$。对任意 $F_A$-代数 $(X,[r_0,r_1])$，唯一的代数同态 $\fold:\List(A)\to X$ 满足结构递归方程：
\begin{align}
\fold(\mathrm{Nil}) &= r_0(\bullet),\\
\fold(\mathrm{Cons}(a,\ell)) &= r_1(a,\fold(\ell)).
\end{align}
\end{proposition}

这两个方程精确地刻画了"折叠"的语义——它们是从代数同态的交换图中逐元素解包得到的。

\begin{remark}[与等变性的类比]
方程 (1) 和 (2) 可以视为\textbf{在列表结构上的广义等变约束}。群等变性说"先作用再映射等于先映射再作用"；列表折叠方程说"先构造列表再折叠等于先折叠子结构再组合"。两者都描述了函数与结构之间的\textbf{协调性}——这正是范畴论中"同态"统一概念的核心。
\end{remark}

\subsection{对偶：余代数同态与结构余递归}

\begin{example}[Mealy 机器的展开]\label{ex:mealy-unfold}
设 $H_{O,I}(X)=I\to O\times X$。从余代数 $(X,n)$ 到 $(\Mealy_{O,I},\mathrm{next})$ 的余代数同态 $\unfold:X\to\Mealy_{O,I}$ 满足：
\begin{align}
\mathrm{next}(\unfold(x))(i)_1 &= n(x)(i)_1,\\
\unfold(\mathrm{next}(\unfold(x))(i)_2) &= \mathrm{next}(\unfold(x))(i)_2.
\end{align}
这里下标 $(\cdot)_1$ 和 $(\cdot)_2$ 分别表示输出分量和下一状态分量。
\end{example}

\textbf{余代数同态的关键洞察：} 它捕获了"展开"（unfold）的合法性条件——展开操作必须与余代数的结构映射协调。这为后续在 Para 范畴中讨论 RNN 的权重共享正确性提供了理论基础。

\newpage

% ==================================================================
\section{从自函子代数到自由单子代数}
% ==================================================================

\subsection{原文翻译与补充}

论文在 2.2 节结尾（注 2.13）提到：
\begin{quote}
\itshape
The fact that these are initial arises from a deeper fact related to the fact that, in many cases, for a given endofunctor there is a monad whose category of monad algebras is equivalent to the original category of endofunctor algebras. We note this because the construction which takes us from one to the other, the so-called algebraically free monad on an endofunctor, will be seen in Part 3 to derive RNNs and other similar architectures from first principles.
\end{quote}

\textbf{翻译：}

列表是始代数这一事实源于一个更深刻的事实：在很多情况下，对给定的自函子，存在一个单子，其单子代数范畴等价于原来的自函子代数范畴。我们指出这一点，是因为从自函子到该单子的构造——即所谓的"自函子上的代数自由单子"——将在第三部分中用于从第一性原理推导 RNN 和其他类似架构。

\subsection{严谨补充：代数自由单子定理}

下面给出完整陈述和证明概要。

\begin{theorem}[代数自由单子的存在性]\label{thm:free-monad}
设 $\cat C$ 是完备且余完备的范畴，$F:\cat C\to\cat C$ 为自函子。若 $F$ 保持 $\kappa$-滤过余极限（对某个正则基数 $\kappa$），则遗忘函子 $U:\Alg_{\Monad}(\FreeMnd(F))\to\Alg_{\Endo}(F)$ 是范畴等价。特别地，$\FreeMnd(F)$ 的代数范畴等价于 $F$ 的代数范畴：
\[\Alg_{\Monad}(\FreeMnd(F))\simeq\Alg_{\Endo}(F).\]
\end{theorem}

\begin{proof}[证明概要]
自由单子 $\FreeMnd(F)=T$ 的构造如构造 \ref{cons:free-monad} 所述。需要验证：
\begin{enumerate}[nosep]
    \item $T$ 上的单子结构由 $F$ 的迭代自然诱导——单位来自 $F^0\hookrightarrow T$ 的嵌入，乘法来自相邻迭代层的拼接；
    \item 从 $F$-代数 $(A,a:F(A)\to A)$ 构造 $T$-代数结构 $\bar a:T(A)\to A$ 需要沿超穷归纳定义，在每个后继步应用 $a$，在极限步取余极限；
    \item 反过来，$T$-代数的结构限制到 $F(A)\hookrightarrow T(A)$ 上就给出 $F$-代数结构。
\end{enumerate}
两个方向互逆，从而得到范畴等价。完整证明见 Kelly (1980) 或 Barr \& Wells (1985)。
\end{proof}

\begin{example}[从 $1+A\times-$ 到列表单子]
对于 $F_A(X)=1+A\times X$，$\FreeMnd(F_A)$ 的承载函子是 $X\mapsto\List(A)\times X$，单子结构为：
\begin{itemize}[nosep]
    \item $\eta_X:X\to\List(A)\times X$，$x\mapsto([\,],x)$；
    \item $\mu_X:\List(A)\times\List(A)\times X\to\List(A)\times X$，$(\ell_1,\ell_2,x)\mapsto(\ell_1{+\!+}\ell_2,x)$（列表拼接）。
\end{itemize}
\end{example}

\textbf{这个定理对深度学习的意义：} 当我们在第三部分将自函子提升到 2-范畴 Para 中时，同样的迭代构造将精确地生成\textbf{循环神经网络的超穷展开}——即 RNN 沿着输入序列逐步展开其隐藏状态的过程。权重共享（weight tying）则对应于展开过程中参数的\textbf{余交换余幺半群结构}（即复制和删除参数的能力）。

\newpage

% ==================================================================
\section{2-范畴与参数化态射 $\Para$}
% ==================================================================

\subsection{引入 2-范畴的动机}

\begin{quote}
\itshape
While category theory is a powerful framework, it leaves much to be desired in terms of higher-order relationships between morphisms. It only deals with sets of morphisms, with no possible way to compare elements of these sets. This is where the theory of 2-categories comes in, which deals with an entire category of morphisms.
\end{quote}

\textbf{翻译：}

范畴论虽然是一个强大的框架，但在态射之间的高阶关系方面还有很多不足。它只能处理态射的集合，无法比较这些集合中的元素。这就是 2-范畴理论的用武之地——它处理的是整个态射范畴。

\textbf{补充：为什么需要 2-范畴？}

在之前的讨论中，我们将神经网络层视为线性等变映射，并用代数同态来描述约束。但这有两个局限：
\begin{enumerate}[nosep]
    \item 大多数神经网络不只是一层——它们是\textbf{参数化}的、由多层复合而成的非线性函数；
    \item 权重共享（同一个参数在不同位置被使用）需要显式地建模\textbf{参数的复制}。
\end{enumerate}
2-范畴提供了"态射之间的态射"（2-态射），这使得我们可以正式地讨论\textbf{重新参数化}——即在不改变函数行为的前提下改变参数表示，包括通过复制映射 $\Delta_P:P\to P\times P$ 来"共享"权重。

\subsection{2-范畴的定义}

\begin{definition}[严格 2-范畴]\label{def:2-category}
一个\textbf{严格 2-范畴} $\cat K$ 由以下数据组成：
\begin{itemize}[nosep]
    \item 一个\textbf{对象}的类（0-态射）；
    \item 对任意 $A,B\in\ob{\cat K}$，一个\textbf{态射范畴} $\cat K(A,B)$——其对象称为\textbf{1-态射}，其态射称为\textbf{2-态射}；
    \item 水平复合函子 $\circ:\cat K(B,C)\times\cat K(A,B)\to\cat K(A,C)$；
    \item 恒等 1-态射 $\id_A\in\cat K(A,A)$。
\end{itemize}
满足严格的结合律和单位律。2-态射通常表示为：
\[
\begin{tikzcd}
A \arrow[rr,bend left,"f"] \arrow[rr,bend right,"g"'] & \Rightarrow\alpha & B
\end{tikzcd}
\]
即 $\alpha:f\Rightarrow g$ 是 $f$ 和 $g$ 之间的 2-态射。
\end{definition}

\begin{example}[$\Cat$ 作为 2-范畴]
所有（小）范畴、它们之间的函子、以及函子之间的自然变换构成了一个 2-范畴 $\Cat$：
\begin{itemize}[nosep]
    \item 0-态射：范畴；
    \item 1-态射：函子；
    \item 2-态射：自然变换。
\end{itemize}
\end{example}

\subsection{$\Para$ 构造}

\begin{definition}[$\Para_\rhd(\cat C)$]\label{def:para}
设 $(\cat M,\otimes,I)$ 为幺半范畴，$(\cat C,\rhd)$ 为 $\cat M$-actegory（作用范畴）。则 $\Para_\rhd(\cat C)$ 是以下 2-范畴：
\begin{itemize}[nosep]
    \item \textbf{对象}：与 $\cat C$ 相同；
    \item \textbf{1-态射} $X\to Y$：对 $(P,f)$，其中 $P\in\ob{\cat M}$，$f:P\rhd X\to Y$ 是 $\cat C$ 中的态射；
    \item \textbf{2-态射} $(P,f)\Rightarrow(P',f'):X\to Y$：$\cat M$ 中的态射 $r:P'\to P$，使得
    \[
    \begin{tikzcd}
    P\rhd X \arrow[rd,"f"] & \\
    & Y \\
    P'\rhd X \arrow[ru,"f'"'] \arrow[uu,"r\rhd\id_X"']
    \end{tikzcd}
    \]
    交换；
    \item \textbf{复合}：$(Q,g):Y\to Z$ 与 $(P,f):X\to Y$ 的复合为 $(Q\otimes P, h)$，其中
    \[h:(Q\otimes P)\rhd X\xrightarrow{\mu_{Q,P,X}} Q\rhd(P\rhd X)\xrightarrow{Q\rhd f} Q\rhd Y\xrightarrow{g} Z;\]
    \item \textbf{恒等态射}：$(I,\lambda_X^{-1})$，其中 $\lambda_X:I\rhd X\xrightarrow{\cong}X$ 是 actegory 的左单位子。
\end{itemize}
\end{definition}

直观上，$\Para$ 范畴的态射将\textbf{参数空间} $P$ 与\textbf{实际计算} $f:P\rhd X\to Y$ 打包在一起。复合时参数空间通过 $\otimes$ 并置，这恰好对应于神经网络中"每一层有自己的参数矩阵，组合时参数张量积拼接"的行为。

\begin{remark}[参数化映射的弦图表示]
参数映射 $(P,f):X\to Y$ 可以用弦图（string diagram）直观表示：参数线垂直绘制，表示参数是态射数据的一部分而非域/余域对象。重新参数化 $r:P'\to P$ 在弦图中垂直接入。由于弦图涉及三维拓扑，这里省略图解，读者可参考论文原文的 string diagram 插图。
\end{remark}

\subsection{2-范畴中的"松散"代数}

在 1-范畴中，一个方块图（四条边组成的交换图）要么交换，要么不交换。但在 2-范畴中，方块图的四条边都是 1-态射，而方块内部可以有一个 2-态射（用双线箭头表示），这个 2-态射从"绕上边再走右边"指向"走左边再走下边"。根据这个 2-态射的性质，我们有四种可能：
\begin{itemize}[nosep]
    \item \textbf{严格交换}（strict）：方块严格交换，没有 2-态射填充；
    \item \textbf{伪交换}（pseudo）：方块被一个可逆的 2-态射（同构）填充，标记 $\cong$；
    \item \textbf{lax 交换}：方块被一个 2-态射 $\Rightarrow$ 从复合的"左-下"方向指向"上-右"方向填充；
    \item \textbf{oplax 交换}：方块被一个 2-态射 $\Leftarrow$ 从"上-右"方向指向"左-下"方向填充。
\end{itemize}

论文的核心主张是：\textbf{lax 代数}已经足以从第一性原理推导出递归、循环和类似神经网络。

\begin{definition}[2-单子的 Lax 代数]\label{def:lax-alg}
设 $(T,\eta,\mu)$ 为 2-范畴 $\cat K$ 上的 2-单子。一个\textbf{lax $T$-代数}是一个四元组 $(A,a,\epsilon_A,\delta_A)$，其中
\begin{itemize}[nosep]
    \item $A\in\ob{\cat K}$ 是承载对象；
    \item $a:T(A)\to A$ 是 1-态射（结构映射）；
    \item $\epsilon_A:\id_A\Rightarrow a\circ\eta_A$ 是 2-态射（lax 单位），表示"恒等几乎等于先应用单位再应用结构映射"；
    \item $\delta_A:a\circ T(a)\Rightarrow a\circ\mu_A$ 是 2-态射（lax 结合律），表示"结构映射几乎满足结合律"。
\end{itemize}
满足 lax 单位律和 lax 结合律的相容条件（详见论文附录 F.2）。
\end{definition}

\subsection{核心定理：Lax 代数产生余幺半群}

\begin{theorem}[定理 G.10]\label{thm:comonoid}
设 $(\cat C,\rhd)$ 为 $\cat M$-actegory，$T:\cat C\to\cat C$ 为带强度的自函子单子。考虑提升后的 2-单子 $\Para_\rhd(T)$ 的 lax 代数 $(A,(P,a),\epsilon_A,\delta_A)$。则 $P$ 是 $\cat M$ 中的\textbf{余幺半群}（comonoid），其中
\begin{itemize}[nosep]
    \item $\epsilon_A$ 提供余单位 $!_P:P\to I$；
    \item $\delta_A$ 提供余乘法 $\Delta_P:P\to P\otimes P$；
\end{itemize}
且余幺半群律由 lax 代数的相容条件精确保证。
\end{theorem}

\textbf{这个定理的意义：} 在 $\Set$（或 $\Vect$）中，余幺半群结构恰好就是\textbf{复制}（$\Delta_P$）和\textbf{删除}（$!_P$）参数的能力。这正是权重共享的范畴论本质——当我们在展开 RNN 时，同一个参数 $P$ 被复制到每个时间步使用，这对应了 lax 2-态射 $\delta_A$ 提供的复制映射。

\newpage

% ==================================================================
\section{参数化 2-自函子与神经网络架构}
% ==================================================================

\subsection{从自函子到参数化 2-自函子}

论文附录 I 展示了如何将第 2 节中的自函子通过 $\Para$ 提升为 2-自函子，其代数/余代数则自然地对应于各种神经网络架构。

\begin{definition}[$\Para(F)$ 的构造]\label{def:para-endo}
设 $F:\cat C\to\cat C$ 是带 actegorical strength $\sigma_{P,X}:P\rhd F(X)\to F(P\rhd X)$ 的自函子。则 $\Para_\rhd(F):\Para_\rhd(\cat C)\to\Para_\rhd(\cat C)$ 是如下 2-自函子：
\begin{itemize}[nosep]
    \item 对象上：与 $F$ 相同；
    \item 1-态射上：将 $(P,f):X\to Y$ 映为 $(P,f')$，其中
    \[f':P\rhd F(X)\xrightarrow{\sigma_{P,X}} F(P\rhd X)\xrightarrow{F(f)} F(Y);\]
    \item 2-态射上：将 $r:P'\to P$ 映为自身。
\end{itemize}
\end{definition}

\subsection{折叠 RNN}

\begin{example}[折叠 RNN 作为 $\Para(1+A\times-)$ 的代数]\label{ex:folding-rnn}
取 $F_A(X)=1+A\times X$。$\Para(F_A)$ 的一个代数由承载集 $S$（隐藏状态空间）和参数化映射
\[(P,\cell):1+A\times S\longrightarrow S\]
组成。利用同构 $P\rhd(1+A\times S)\cong P+(P\times A\times S)$（在 $\Set$ 和 $\Vect$ 中），$\cell$ 分解为：
\begin{itemize}[nosep]
    \item $\cell_0:P\to S$（初始隐藏状态，由参数生成）；
    \item $\cell_1:P\times A\times S\to S$（RNN 单元的更新函数）。
\end{itemize}

通过第 4 节讲的自由单子构造（现在在 $\Para$ 中执行），我们得到\textbf{展开后的折叠 RNN}——它沿着输入列表逐步应用 $\cell_1$，而权重共享由 $\Delta_P:P\to P\times P$（复制参数）保证。
\end{example}

\subsection{展开 RNN}

\begin{example}[展开 RNN 作为 $\Para(O\times-)$ 的余代数]\label{ex:unfolding-rnn}
取 $F_O(X)=O\times X$。$\Para(F_O)$ 的一个余代数由承载集 $S$ 和参数化映射
\[(P,\langle\cell_o,\cell_n\rangle):S\longrightarrow O\times S\]
组成，其中 $\cell_o:P\times S\to O$ 计算输出，$\cell_n:P\times S\to S$ 计算下一状态。

通过终余代数的万有性质（类似于 $\Stream(O)$ 是 $O\times-$ 的终余代数），我们得到展开后的 RNN——从初始状态 $s_0$ 出发，无限地产出流 $\Stream(O)$：
\[s_0\xrightarrow{\cell_n}s_1\xrightarrow{\cell_n}s_2\xrightarrow{\cell_n}\cdots\]
每一步输出由 $\cell_o(p,s_i)$ 给出，而参数 $p$ 每一步都被复制使用。
\end{example}

\subsection{递归 NN}

\begin{example}[递归 NN 作为 $\Para(A+(-)^2)$ 的代数]\label{ex:recursive-nn}
取 $G_A(X)=A+X^2$。$\Para(G_A)$ 的一个代数给出\textbf{树递归神经网络}：承载集 $S$，参数化映射
\[(P,\cell):A+S^2\longrightarrow S,\]
分解为叶节点处理 $\cell_0:P\times A\to S$ 和内部节点组合 $\cell_1:P\times S^2\to S$。展开后对二叉树自底向上递归计算，参数每一步都被共享。
\end{example}

\subsection{Mealy 和 Moore 机器作为参数化余代数}

\begin{example}[全 RNN 作为 $\Para(I\to O\times-)$ 的余代数]
取 $H_{O,I}(X)=I\to O\times X$。$\Para(H_{O,I})$ 的一个余代数精确对应对\textbf{可学习的 Mealy 机器}——即全循环神经网络（full RNN）。其核心单元是一个参数化映射
\[n:P\times S\times I\longrightarrow O\times S,\]
在每个时间步接收隐藏状态 $S$、输入 $I$，产生输出 $O$ 和新的隐藏状态 $S$。参数 $P$ 在所有时间步共享。
\end{example}

\begin{example}[Moore 机器 NN 作为 $\Para(O\times(I\to-))$ 的余代数]
类似地，Moore 机器型的神经网络对应 $\Para(O\times(I\to-))$ 的余代数。与 Mealy 型（全 RNN）的关键区别在于：Moore 型的输出 $O$ 不依赖于当前输入 $I$，只依赖于当前状态。
\end{example}

\newpage

% ==================================================================
\section{神经网络的超穷展开}
% ==================================================================

\subsection{展开作为 Lax 代数同态}

论文附录 J 展示了自动机/神经网络展开的精确 2-范畴公式。核心思想是：一个 RNN 单元的\textbf{展开}（unrolling）是如下 lax 代数同态：

\begin{theorem}[折叠 RNN 的展开，例 J.1]\label{thm:unroll-folding}
设 $(P,\cell)$ 是 $\Para(1+A\times-)$ 的代数（即折叠 RNN 单元）。其展开是如下交换图（带 lax 2-态射 $\Delta_P$）：
\[
\begin{tikzcd}
1+A\times\List(A) \arrow[r,"\id+\id_A\times\fold"] \arrow[d,"{[\mathrm{Nil},\mathrm{Cons}]}"'] & 1+A\times S \arrow[d,"{(P,\cell)}"] \\
\List(A) \arrow[r,"{(P,\fold)}"'] & S
\end{tikzcd}
\qquad
\text{带 2-态射 } \Delta_P: P \to P\times P
\]
其中 $\fold:P\times\List(A)\to S$ 是结构递归函数：
\begin{align*}
\fold(p,\mathrm{Nil}) &= \cell_0(p),\\
\fold(p,\mathrm{Cons}(a,as)) &= \cell_1(p,a,\fold(p,as)).
\end{align*}
$2$-态射 $\Delta_P:P\to P\times P$ 是参数复制映射，它保证了"tying"——每一步展开使用相同的参数 $p$。
\end{theorem}

\textbf{交换条件验证：} 展开图的上右复合得到 $(P\times P, \text{上右})$，下左复合得到 $(P, \text{下左})$。$\Delta_P$ 是它们之间的 2-态射，即重新参数化——因为下左的函数在每一步都使用同一个 $p$，而上右的函数在第一个参数和递归参数之间使用两个独立的 $p,q$。$\Delta_P$ 将 $p$ 复制为 $(p,p)$，强制它们相同——这正是\textbf{权重共享}的范畴论表述。

\subsection{展开 RNN 的展开}

\begin{theorem}[展开 RNN 的展开，例 J.2]
设 $(P,\langle\cell_o,\cell_n\rangle)$ 是 $\Para(O\times-)$ 的余代数。其展开由以下结构递归给出：
\[\unfold(p,s)=\mathrm{MkStream}(\cell_o(p,s),\;\unfold(p,\cell_n(p,s))),\]
其中权重共享由参数复制的 lax 2-态射 $\Delta_P$ 保证。
\end{theorem}

\textbf{关键观察：} 在所有这些展开中，$2$-态射部分始终是 $\Delta_P:P\to P\times P$（或更一般地，$\delta_A$ 给出的余幺半群余乘法）。这统一地解释了为什么不同神经网络架构中的权重共享都可以用同一个数学机制来描述。

\subsection{统一视角}

所有架构的核心——无论是折叠 RNN、展开 RNN、递归 NN、Mealy 机器还是 Moore 机器——都可以总结为同一个模板：

\begin{center}
\begin{tabular}{c|c|c|c}
\toprule
\textbf{架构} & \textbf{自函子} & \textbf{代数/余代数} & \textbf{展开结果} \\
\midrule
折叠 RNN & $1+A\times-$ & 代数 & 折叠（结构递归） \\
展开 RNN & $O\times-$ & 余代数 & 展开/流（结构余递归） \\
递归 NN & $A+(-)^2$ & 代数 & 树折叠 \\
全 RNN（Mealy） & $I\to O\times-$ & 余代数 & Mealy 机器展开 \\
Moore RNN & $O\times(I\to-)$ & 余代数 & Moore 机器展开 \\
\bottomrule
\end{tabular}
\end{center}

而\textbf{权重共享}的实质是：在 $\Para$ 范畴中，lax 代数/余代数的 lax 2-态射提供了参数空间上的\textbf{余幺半群结构}（复制和删除），这在展开的超穷构造中每一步都被施加，确保同一参数被反复使用。

\newpage

% ==================================================================
\section{总结与展望}
% ==================================================================

\subsection{论文核心贡献总结}

\begin{enumerate}[nosep]
    \item \textbf{GDL 的范畴推广：} 几何深度学习中的等变性 = 群作用单子的代数同态。这统一了所有对称性约束。
    \item \textbf{自函子代数统一数据结构：} 列表、树、自动机都可以视为特定自函子的（余）代数。折叠/展开操作是（余）代数同态。
    \item \textbf{2-范畴 $\Para$ 建模参数化：} $\Para(\cat C)$ 的 1-态射将参数空间显式打包，2-态射捕获重新参数化（包括权重共享）。
    \item \textbf{权重共享 = 余幺半群结构：} $\Para$ 中 lax 代数的 lax 2-态射自动产生参数空间上的余幺半群（comonoid）结构，即复制和删除参数的能力。
    \item \textbf{架构展开 = 自由单子构造：} 循环/递归网络的展开对应于自函子的超穷迭代（自由单子构造），这是始代数/终余代数万有性质的 2-范畴版本。
\end{enumerate}

\subsection{开放问题与未来方向}

论文第 4 节提出了若干开放方向，这里择要列出：

\begin{itemize}[nosep]
    \item \textbf{非线性处理：} GDL 通常只能处理单个线性层的等变约束。CDL 框架通过 lax 代数可以跨越多个层和非线性激活函数来推理。
    \item \textbf{类型安全的神经网络：} 选择合适的多项式函子作为自函子，可以编码类型约束，使得神经网络只能学习\textbf{良类型}的函数。
    \item \textbf{可验证的 AI：} 通过范畴逻辑，可以精确指定架构能使用的推理形式，这为实现可验证公平性和安全性提供了数学基础。
    \item \textbf{函数式微分：} 结合 Nunes \& Vákár (2023) 关于归纳/余归纳类型的可微性结果，可以在纯函数式语言中实现端到端可微的架构设计。
\end{itemize}

\vspace{1em}
\begin{center}
\rule{0.5\textwidth}{0.5pt}\\[0.5em]
\textbf{教程完}
\end{center}

\end{document}